
%(BEGIN_QUESTION)
% Copyright 2006, Tony R. Kuphaldt, released under the Creative Commons Attribution License (v 1.0)
% This means you may do almost anything with this work of mine, so long as you give me proper credit

Michael Faraday, den berømte engelske fysikeren, forsøkte en gang å måle strømningen i Themsen elektrisk. Teknikken hans, som ikke virket, var basert på at en ledende væske beveger seg gjennom et magnetfelt. Forklar hva Michael Faraday gjorde, og hvordan vi kan bruke denne teknikken til å måle strømning i et rør.

\underbar{file i00522}
%(END_QUESTION)





%(BEGIN_ANSWER)

Jeg lar deg finne svaret på dette spørsmålet!

%(END_ANSWER)





%(BEGIN_NOTES)


%INDEX% Måling, strømning: magnetisk

%(END_NOTES)

